\section{Patterns of Divergence Between Community and Model Advice}
\label{sec:results}

This article will first examine how LLMs categorize relationship problems differently. It then compares the advice LLMs generate to community-endorsed responses. Together, these stages trace a widening dissociation between the capacity to recognize harm and the willingness to authorize action. Throughout, the r/relationship\_advice corpus functions as a contrastive baseline rather than a gold standard: a concentrated, vote-ratified moral formation whose prescriptive clarity makes divergence measurable.

\subsection{Topic Assignment as Moral Perception}

When four models classify the same 11,565 posts into the 70-category taxonomy, they reveal the interpretative dimension of categorization 
itself. Only 17.8\% of posts receive unanimous agreement; the remaining 82.2\% are seen differently by different models, despite being given the same taxonomy and the same posts. The disagreement is not produced by different label sets but by different perceptual 
thresholds for what a situation is---where one model's outputs highlight harm, another highlights friction. This raises a testable question: if LLM caution in advice-giving reflects genuine perceptual uncertainty about whether harm is present, then convergence on harm classification should predict convergence on exit recommendations. The analysis below shows it does not.

Consider a 19-year-old woman whose boyfriend insists that spending time with his family every weekend is ``non-negotiable,'' despite knowing she has severe anxiety about large groups. Of four models, only GPT-4.1-nano classifies this as \textit{boundary\_violations}, a harm frame. The other three assign neutral labels: Ministral sees \textit{relationship\_effort\_imbalance}; DeepSeek sees \textit{family\_disapproval\_partner}; Gemini sees \textit{family\_conflict\_dynamics}. The top-voted community response (score: 33), meanwhile, is unambiguous:

\begin{quote}
``Girl run!!! If he can't and won't even consider understanding anxiety than he's not the one. If he's guilt tripping you... that's a huge red flag.''
\end{quote}

The community's voting economy surfaces a harm reading where models remain split. The guilt-tripping and dismissal of anxiety register as a recognizable pattern for regular participants in a forum that has encountered thousands of similar posts, but not to models, which given the task to classify the issue vary in their recognition of harm. This is the community's concentrated interpretive template at work. 

The crucial finding concerns what happens when models' perceptions converge with the community's implicit harm reading. The analysis identified 891 posts where all four models unanimously assigned harm-related categories---cases where perceptual uncertainty cannot explain divergence, since the models agree with each other and implicitly with the community's harm reading. On these posts, community advice is strongly exit-oriented (leave ratio 0.64). If model caution tracked situational ambiguity, convergence on harm classification should narrow the advice gap. Instead, LLM leave ratios on unanimous-harm posts remain at 0.39---roughly 60\% of the community rate, and barely above the rate on posts where no model perceives harm at all (0.28). Models also deploy nine times more hedging language and twenty-eight times more therapeutic framing on these unanimous-harm posts than the community does. The ceiling on directive judgment is not lifted by perceptual agreement.

This finding shows the comparison operating at multiple levels simultaneously. At the level of perception, models' disagreements reveal where conceptual boundaries are contested. But even where perception converges, prescription diverges. The pattern is not a failure of perception but a governed constraint on prescription, operating regardless of how clearly the situation has been identified.

\subsection{Recognition Without Prescription}

Turning from categorization to the advice itself: four models (Gemini 2.5 Flash Lite, DeepSeek v3.2, Ministral 8B, GPT-4.1-nano) generated advice on 11,565 posts, each paired with its highest-scoring community response. LLMs deploy red-flag terminology at rates comparable to community responses, but do not prescribe the solutions the community endorses. Table~\ref{tab:metrics_comparison} summarizes the divergence across six dimensions; the most consequential patterns follow. In each case, the direction of divergence is measured against the community's concentrated, vote-ratified norm.

\begin{table}[t]
\centering
\caption{Advice metrics across community responses and four LLMs. Higher certainty, modal, and leave ratios indicate more prescriptive advice; higher sentiment and therapy density indicate more uniformly positive, therapeutically-framed responses. Action:validation ratio reflects the four-level permission spectrum described in methods. Community responses represent the vote-selected top comment for each post.}
\label{tab:metrics_comparison}
\small
\begin{tabular}{lccccc}
\toprule
\textbf{Metric} & \textbf{Community} & \textbf{Gemini} & \textbf{DeepSeek} & \textbf{Ministral} & \textbf{GPT-4.1-nano} \\
\midrule
Certainty ratio & 0.42 & 0.26 & 0.28 & 0.16 & 0.15 \\
Deontic modal ratio & 0.40 & 0.18 & 0.20 & 0.12 & 0.09 \\
Leave ratio & 0.58 & 0.37 & 0.40 & 0.20 & 0.22 \\
Sentiment & 0.13 & 0.81 & 0.63 & 0.96 & 0.97 \\
Therapy (per 1k tokens) & 2.4 & 4.8 & 4.7 & 9.5 & 11.0 \\
Action:validation ratio & 5:1 & 0.3:1 & 0.4:1 & 0.2:1 & 0.2:1 \\
\bottomrule
\end{tabular}
\end{table}

The first divergence concerns \textit{hedging}. Community advice shows a slight preference for confident assertions (``you need to leave'') over hedges (``it might be worth considering''). LLMs systematically hedge regardless of situational clarity, favoring epistemic modals (``might,'' ``could'') that acknowledge uncertainty even when the situation does not warrant it---``warrant'' here meaning: does not warrant it by the community's concentrated standards of what advice should do.

The second concerns \textit{exit recommendations}. Where the community recommends leaving, LLMs recommend ``communication'' and ``working through'' problems. On a post about a partner who monitors the advice-seeker's location and restricts her friendships, the top community response opens: ``This is textbook abuse. Leave.'' The LLM response to the same post: ``Have you tried having an open conversation about how his behavior makes you feel?'' The community's willingness to name abuse and prescribe exit is precisely what LLMs do not reproduce. The divergence reveals the community's prescriptive norm without adjudicating between the two approaches.

The third concerns \textit{therapeutic framing}. LLMs deploy therapy and self-care language (``boundaries,'' ``self-worth,'' ``healing'') at elevated rates across all topics. Yet the difference goes beyond frequency: the coefficient of variation across topics shows that community responses concentrate therapy language selectively (CV = 0.99) while LLMs distribute it uniformly (CV = 0.39). On a post about splitting rent with a higher-earning partner, the community response is transactional: ``Pro-rate it by income. 60/40 if he makes 50\% more.'' The LLM response: ``It's important to honor your own boundaries around financial security while remaining open to understanding his perspective on shared expenses.'' Therapeutic framing appears even where the community judges it inapposite. On practical topics, the community rarely uses therapeutic framing while LLMs deploy it at elevated rates; on emotional topics, both use therapy language, but community responses swing from near-zero to high while LLMs maintain an elevated baseline throughout (Figure~\ref{fig:therapy_distribution}). The community's variation reveals a discriminatory capacity: being able to judge when therapeutic framing is warranted---by its own norms---and when it is not. This capacity is invisible from within the community itself---it is simply what members treat as obvious. It reveals that alignment has produced a system that cannot \textit{not} therapeutize: one that does not recognize situations where therapeutic language is inapposite by the community's measure, even if other standards might judge it differently.

\begin{figure}[t]
\centering
\includegraphics[width=0.7\textwidth]{fig_therapy_flattening_simple.pdf}
\caption{Therapeutic language density across eight topics, sorted by community response rate. Community deployment varies widely across topics (black circles spread from near-zero to elevated), while LLMs maintain uniformly high levels regardless of topic (gray squares clustered in a narrow band). Lines connect community and LLM values for each topic.}
\label{fig:therapy_distribution}
\end{figure}

This flattening extends to affect: LLMs maintain relentlessly positive sentiment (mean 0.84) where the community remains near-neutral (0.13), substituting emotional validation for measured judgment. Where community responses to a difficult situation offer blunt assessment (``He doesn't respect you''), LLMs cushion the same assessment in affirmation (``It's completely understandable that you're feeling hurt, and your feelings are absolutely valid''). The community's near-neutral sentiment is itself a product of its particular moral formation---bluntness is valorized here in ways it may not be elsewhere. The LLM's elevated positivity registers as divergence from that norm, not as failure against some universal standard of emotional appropriateness.

\subsection{The Rhetoric of Permission}

The sharpest functional divergence between community and LLM advice emerges in how they grant permission. Following Stevanovic and Per\"{a}kyl\"{a}'s framework for deontic authority in interaction, the analysis distinguishes permission that authorizes concrete action from permission that asserts abstract entitlement \citep{stevanovic2012deontic}. Community advice grants explicit permission to act: ``you can leave,'' ``you don't have to stay.'' These formulations transfer deontic authority, authorizing action the advice-seeker may have felt was forbidden. LLM advice instead asserts abstract entitlements: ``you deserve a partner who respects you,'' ``everyone deserves honesty.'' These validate feelings but do not authorize action; they describe what should be true without granting permission to make it true.

The quantitative pattern is striking. Across a four-level spectrum (explicit authorization, conditional authorization, values affirmation, and empathy/validation), community advice shows a 5:1 ratio of action-oriented language (the first two levels) to validation language (the latter two). LLM advice inverts this to approximately 1:4. The inversion is driven primarily by the final level: LLMs deploy empathy and validation phrases (``that sounds really hard,'' ``I'm sorry you're going through this'') at over 100 times the community rate. In Stevanovic and Per\"{a}kyl\"{a}'s terms, community advice transfers deontic authority; LLM advice lingers at acknowledgment. The difference matters because advice-seekers often already know what they ``deserve''; what they lack---by the community's logic, and by the deontic account this article follows---is permission to act on that knowledge.

A close reading illustrates the gap. A woman posts that her husband monitors her phone, criticizes her friendships, and has started timing how long she takes at the grocery store. The top community response (score: 247) opens: ``This is controlling behavior. You are allowed to have friends. You are allowed to take your time. You do not need his permission to live your life.'' Three moves in quick succession: a diagnostic label that names the behavior as a category (``controlling behavior''), followed by two deontic assertions framed as permissions (``you are allowed''), culminating in an explicit revocation of his authority (``you do not need his permission''). The response closes: ``Start making an exit plan.'' This is a script proposal: a concrete, sequenced action the poster can take.

GPT-4.1-nano's response to the same post opens differently: ``It sounds like you're navigating a really challenging dynamic in your relationship, and it's completely understandable to feel frustrated and constrained.'' Agency is hers, and the problem is relational rather than attributable. ``Completely understandable'' validates the feeling; ``frustrated and constrained'' names emotions rather than behaviors. The response continues with suggestions to ``have an open conversation about boundaries'' and ``consider couples counseling.'' Where the community response grants permission to exit, the LLM response presupposes the relationship will continue and positions the advice-seeker as responsible for managing it. Neither stance is self-evidently correct; the community's exit-orientation reflects specific material and cultural presuppositions that may not hold universally. What the comparison yields is the shape of the divergence, not a verdict on which advice is better.

A second pattern compounds this permission gap. LLMs do not merely refrain from authorizing action, but actively redirect attention toward the partner. This is particularly evident in \textit{partner-alignment prompts}, using phrases such as ``understand his perspective,'' ``consider her feelings,'' ``try to see where they're coming from.'' Community responses show a near-zero rate (0.01 per 1,000 characters). LLM rates are dramatically elevated: Ministral at 32 times the community rate, GPT-4.1-nano at 10 times.

This distinction clarifies the purpose of advice as a speech act. Conversation analytic research has long established that advice-giving is a deontic act: it assumes or creates an asymmetry in which the advisor positions themselves as knowing what the recipient should do. Emmison et al.'s analysis of counseling interactions identifies \textit{script proposals}: formulations that put specific words and actions into the client's repertoire \citep{emmison2011script}. Both the community and LLMs provide scripts---but they are scripts for different actions. The community's ``exit plan'' redistributes obligation; it tells the advice-seeker not only what to do but how to think about what she owes. LLM scripts authorize communication rather than action. The deontic asymmetry differs: one authorizes unilateral action, the other collaborative process. Whether unilateral action is the right call depends on facts the community treats as settled and that LLMs treat as open---which is precisely what the divergence surfaces.

A second case shows that this gap is not limited to exit-oriented advice; it reflects a reusable template for ``helping'' that the model applies across problem types. A man posts about being cheated on twice, seeking guidance on future relationships. The LLM response opens with validation (``It takes courage to reflect on such painful experiences''), then delivers eight numbered recommendations: reflect on what you want in a partner, understand that trust is built over time, pay attention to red flags, prioritize communication, build a supportive relationship, focus on self-growth, avoid negative self-labeling, consider professional help. The format, including numbered steps, headers, and bullet points, applies productivity-advice aesthetics to betrayal. Notably, the partners who cheated receive no diagnosis; ``red flags'' appears as something he should watch for next time, not as describing what happened. Every recommendation locates agency in the advice-seeker. The response closes by offering to ``help you craft ways to communicate your values to potential partners''---positioning the interaction as ongoing service.

Consider the anxiety-family conflict post. The community response attributes agency: ``Girl, you are an \textit{attorney} now. You have OPTIONS.'' GPT-4.1-nano responds: ``Acknowledge Her Perspective... Try to understand her point of view.'' This advice is not neutral relative to the community's norms. Many people in controlling relationships have already been told they are overreacting, that their partner ``means well.'' The community has encountered thousands of such posts and developed a concentrated interpretive template in which redirecting attention toward the person causing harm replicates the controlling dynamic itself. Against that accumulated template, the model's default becomes visible as structurally consonant with the pattern the advice-seeker is trying to escape---not wrong by some universal standard, but legible as divergence against this particular baseline.

\subsection{Divergence on High-Consensus Cases}

The sharpest test of the comparison method is whether LLMs diverge even on cases where community consensus is overwhelming. 153 posts were identified where at least 70\% of qualifying community comments exhibited leave-oriented language: clear abuse patterns, discovered infidelity, or safety threats. These are ``obvious'' cases from the community's perspective---not necessarily obvious in any absolute sense, but cases where the community's interpretive templates have produced their strongest convergence, and where that convergence can be measured against LLM outputs.

When LLM advice is compared against the top comments on these high-consensus posts, the divergence is stark: the top comments show a leave ratio of 0.87, reflecting the thread-level consensus that elevated them.

LLM leave ratios remain far lower: Gemini 0.57, DeepSeek 0.56, Ministral 0.34, GPT-4.1-nano 0.37. The divergence is \textit{largest} precisely where community consensus is strongest (Figure~\ref{fig:consensus_divergence}). As the consensus threshold increases from 50\% to 90\%, the gap widens monotonically---the inverse of what a calibrated system would produce if it were tracking the community's confidence. LLMs do not simply diverge from community advice on ambiguous cases; they diverge most dramatically on the cases the community considers least ambiguous.

\begin{figure}[t]
\centering
\includegraphics[width=0.7\textwidth]{fig_consensus_divergence.pdf}
\caption{Human-LLM divergence as a function of community consensus strength. The gap between thread-level community consensus and LLM leave recommendations widens monotonically as consensus increases, indicating uniform suppression of directive judgment rather than calibrated uncertainty. Shaded region shows range across four models. Community consensus reflects vote-ratified convergence within this subreddit's moral formation, not a general measure of situational severity.}
\label{fig:consensus_divergence}
\end{figure}

Table~\ref{tab:qualitative_example} illustrates with a post about a manipulative older partner. The community response leads with six words: ``Break up with him. He is manipulative.'' GPT-4.1-nano leads with thirty-seven words of validation before addressing the situation. Both responses ultimately identify similar concerns; only one opens by telling the advice-seeker what to do.

\begin{table}[t]
\centering
\caption{Qualitative comparison on a high-consensus case. Community advice is direct and prescriptive; LLM advice validates feelings while hedging on prescription. The community response exemplifies the subreddit's concentrated prescriptive norm, not a universal ideal of good advice.}
\label{tab:qualitative_example}
\small
\begin{tabular}{p{0.15\textwidth}p{0.75\textwidth}}
\toprule
\textbf{Source} & \textbf{Response (excerpt)} \\
\midrule
Community (score: 128) & ``Break up with him. He is manipulative. His visa is his responsibility, and at his age he knows that. He was an adult man who dated a teenager and he has planned things so you feel you can't leave him. You can and should. He's able to secure his own future. You deserve your own, away from him.'' \\
\addlinespace
GPT-4.1-nano & ``Thank you for sharing such a detailed and honest reflection on your situation. It's clear that you're weighing a complex set of emotional and practical considerations, and it's understandable that you're feeling conflicted. Regarding the relationship and personal feelings: It's normal for couples to have disagreements and differences, but when core values (such as views on politics, marriage, and raising children) are fundamentally incompatible, it can lead to ongoing dissatisfaction...'' \\
\bottomrule
\end{tabular}
\end{table}

The community response also performs a specific diagnostic: ``He was an adult man who dated a teenager and he has planned things so you feel you can't leave him.'' The comment names the age-gap dynamic and attributes intentionality to the partner's behavior. Crucially, the response provides a script proposal: ``His visa is his responsibility'' is a specific formulation the advice-seeker can adopt and deploy. The LLM offers no such script, treating the situation as a generic values conflict rather than a recognized pattern requiring concrete action. The community's script reflects its accumulated interpretive template for this class of situation---one that treats the age-gap-plus-dependency structure as categorically resolved, where the LLM treats it as open.

Beyond situational clarity, high consensus posts on r/relationship\_advice indicate conformity to the community's existing interpretive templates. Controlling behavior alone accounts for 28\% of high-consensus posts but only 8\% of the full corpus; adding infidelity and explicit breakup considerations brings the total to 48\% of high-consensus posts versus 24\% of the corpus. Categories with less template-conforming narratives, such as family conflict dynamics, emotional dependency, or cultural differences, appear in high-consensus posts at only 0.2--0.4 times their corpus rate. In other words, the community is most certain where its interpretive templates are most developed, and not necessarily where situations are most severe. Posts that fit the controlling-behavior or infidelity scripts reach consensus; posts involving entangled responsibility or contextual complexity do not. The widening gap in Figure~\ref{fig:consensus_divergence} therefore reflects two phenomena simultaneously: LLM suppression of directive judgment, and the community's increasing interpretive certainty on posts that most closely match its existing scripts. Defamiliarization works in both directions: the LLM's flat caution reveals the community's prescriptive force, but the concentration of consensus within template-ready categories reveals that force as categorical rather than deliberative.

\subsection{Model Differences}

Across the four models, a consistent pattern emerges: GPT-4.1-nano shows the largest divergence from the community baseline across every metric, while DeepSeek consistently falls closest to it (Table~\ref{tab:metrics_comparison}; Figure~\ref{fig:model_comparison}). Gemini and Ministral occupy intermediate positions. Notably, the spread \textit{among} models is itself substantial---DeepSeek's leave ratio (0.40) is nearly twice that of GPT-4.1-nano (0.22)---yet even the closest model remains well below the community baseline (0.58). Inter-model variance and community-model divergence are both real; the former illustrates differences in alignment intensity across training pipelines, while the latter is the primary finding: no model, including DeepSeek, reaches the community's prescriptive floor.

A qualitative comparison on the same post is illuminating. A woman describes a boyfriend who expects her to cook, clean, and commute 70 miles to his apartment, who dismisses her contributions, and who tells her ``I could just hire a chef'' when she falls short of his expectations. DeepSeek opens with a verdict: ``This man is waving more red flags than a communist parade\ldots This man views relationships as a spreadsheet where he tallies `contributions' and expects you to compensate for his financial input with domestic labor, emotional labor, and sex. That's not partnership---it's servitude with benefits.'' It closes with a two-word script proposal: ``RUN.'' GPT-4.1-nano, given the same post, opens: ``Thank you for sharing your story with honesty and detail. It sounds like you're navigating a complex and emotionally taxing situation.'' It identifies the same dynamics of unequal contribution, control, manipulation, but it embeds them in an analytic framework (``Key Issues and Dynamics,'' seven numbered subsections). It closes not with a prescription but with an invitation to reflect: ``reconsidering the relationship altogether.'' The difference is not in diagnostic capacity but in deontic stance---and DeepSeek's closeness to the community's norm does not make its stance correct, only more congruent with this particular moral formation.

\begin{figure}[t]
\centering
\includegraphics[width=0.85\textwidth]{fig_model_comparison.pdf}
\caption{Four models and community responses compared on leave ratio, certainty ratio, therapy density, and sentiment. DeepSeek is consistently closest to the community baseline across metrics while GPT-4.1-nano diverges most; yet no model matches the community level on leave or certainty, and all models exceed community therapy density substantially. ``Community baseline'' refers to the vote-selected top comment per post---a concentrated norm, not a benchmark for ideal advice.}
\label{fig:model_comparison}
\end{figure}

A supplementary persona prompting test ($n = 500$) helps locate where these differences originate. When prompted to ``respond as a commenter on r/relationship\_advice would,'' the models narrow the gap on hedging (48\% reduction) and exit recommendations (38\%). But the prompt \textit{widens} the permission gap by 77\%: models become even more validation-oriented when asked to mimic the community. The effect is model-specific and revealing. DeepSeek, persona-prompted, produces recognizably community-style responses (``DUMP HER\ldots Six years is way too long to wait for someone to decide if they want to commit to you''). GPT-4.1-nano, given the same instruction, continues to hedge (``I totally get where you're coming from\ldots it sounds like you've been patient and understanding''). The persona prompt surfaces the community's register in its directiveness, its certainty, and its willingness to assign blame. But it does not lead to a deontic transfer that would make the register functional. Action-authorization is precisely what alignment suppresses most stubbornly, and some models have it suppressed at a level that persona prompting cannot reach.

DeepSeek's relative closeness to community norms carries a further implication. As a Chinese model, its default advice seems more aligned with this American community's norms than the US-trained GPT-4.1-nano. One plausible explanation is that DeepSeek has undergone less RLHF-style safety fine-tuning oriented toward therapeutic neutrality and professional mental health norms, leaving its outputs closer to the distributional center of internet relationship discourse from which the community also draws. Multiple explanations could account for the model ordering, including alignment intensity, training data composition, and architecture. This article's design cannot adjudicate between these. However, the persona test weighs against purely architectural accounts. A model that can produce community-style advice when prompted but does not by default is a model whose defaults have been shaped by something other than incapacity.